% ============================================================
% IEEE Access / IEEE Transactions on Affective Computing
% Tri-Modal Depression Risk Detection
% Author: Md. Mursalin, Netrokona University
% ============================================================

\documentclass[journal]{IEEEtran}

% ---- Packages ----
\usepackage{amsmath,amssymb,amsfonts}
\usepackage{algorithm}
\usepackage{algorithmic}
\usepackage{graphicx}
\usepackage{textcomp}
\usepackage{xcolor}
\usepackage{booktabs}
\usepackage{multirow}
\usepackage{url}
\usepackage{cite}
\usepackage[hidelinks]{hyperref}

% In Overleaf: upload all fig*.png to project root
\graphicspath{{./}}

% ---- Title ----
\title{Multi-Task Learning for Tri-Modal Depression Detection:
A Controlled Ablation Study on DAIC-WOZ}

\author{Md.~Mursalin
and~Md.~Shovon%
\thanks{Manuscript received \today.
Md. Mursalin is with the Department of Computer Science and
Engineering, Netrokona University, Netrokona 2400, Bangladesh
(corresponding author, e-mail: mursalinshuvo27@gmail.com).
Md. Shovon, Department of Computer Science and Engineering, Netrokona
University, supervised this work.
Dataset: DAIC-WOZ provided by USC Institute for Creative Technologies.
Code: \url{https://github.com/DevMursLab/Thesis}.}%
}

\markboth{IEEE Access,~Vol.~XX, No.~XX, 2025}%
{Mursalin \MakeLowercase{\textit{et al.}}: Multi-Task Tri-Modal Depression Detection}

\begin{document}

\maketitle

% ============================================================
\begin{abstract}
Automated depression screening from clinical interviews requires
multimodal analysis, demographic fairness, and clinical
interpretability, yet prior work rarely evaluates all three under
controlled conditions. We present a tri-modal (facial action units,
acoustic features, transcript text) depression-risk framework on the
DAIC-WOZ corpus, using a single \texttt{ConfigurableFusionModel}
class across seven modality subsets (5~seeds each) to isolate genuine
modality contribution from architectural confounds. Our most
promising observed result is symptom-level \emph{multi-task
learning}---jointly predicting the binary label, PHQ-8 score, and
eight symptom items---which raises mean macro-F1 from $0.607\pm0.063$
to $0.629\pm0.021$ on the official dev split ($N{=}34$) with markedly
reduced seed variance; however, a paired Wilcoxon test matching each
multi-task seed against the single-task baseline run from the same
seed does not reach significance ($p{=}0.625$, multi-task wins
3/5~pairs), so we report this as an encouraging but statistically
unconfirmed direction rather than a validated headline result. On
the held-out test split ($N{=}47$), the
model achieves F1\,$=0.585$ and AUC\,$=0.706$, exceeding the dev AUC
and indicating generalization beyond the dev set; male and female
TPR are numerically equal on this split (0.286 each), though both
are low in absolute terms and pooled 10-fold cross-validation over
all 188 participants shows this equality is not stable across folds
(mean TPR gap~$0.415\pm0.211$). We additionally
evaluate cross-modal attention fusion and an Equalized Odds fairness
loss in 10-seed paired Wilcoxon experiments; neither reaches
significance at this sample size (F1 $p{=}0.496$; TPR-gap
$p{=}0.813$), and post-hoc power analysis shows the study is
genuinely underpowered for these effect sizes rather than showing
true null effects. We report both as transparent negative/inconclusive
findings rather than reframing them as successes. Three convergent
explainability methods (attention rollout, gradient$\times$input
saliency, occlusion ablation) independently identify audio and AU04
(Brow Lowerer) as the primary predictive cues, consistent with
established clinical markers of depressed affect. Our contribution is
methodological rigor and transparent reporting---including honest
limitations---rather than state-of-the-art performance.
\end{abstract}

\begin{IEEEkeywords}
Depression detection, multi-task learning, PHQ-8, multimodal fusion,
DAIC-WOZ, CLNF action units, cross-modal attention, fairness-aware
training, Equalized Odds, explainable AI.
\end{IEEEkeywords}

% ============================================================
\section{Introduction}
\label{sec:intro}

\IEEEPARstart{D}{epression} affects an estimated 280 million people
worldwide \cite{who2023} and is a leading cause of disability globally.
Clinical diagnosis relies on structured interviews, self-report
questionnaires (PHQ-8, PHQ-9), and clinician observation---a
resource-intensive process that is inaccessible in many healthcare
settings. Automated screening from clinical interview recordings offers
a scalable complement to traditional pathways, provided the system
meets three requirements: (1)~\emph{multimodal coverage} matching the
multi-channel nature of clinical observation; (2)~\emph{demographic
fairness} ensuring equal detection rates across gender groups; and
(3)~\emph{clinical interpretability} enabling a human clinician to
audit the model's reasoning.

Prior work on the DAIC-WOZ corpus \cite{gratch2014} has explored
audio features \cite{ringeval2017}, text-based models, and multimodal
fusion \cite{gong2017,williamson2016}. However, three critical gaps
persist. First, ablation studies in existing work typically compare
configurations with \emph{different} architectural choices per modality
subset, confounding the modality contribution with incidental model
changes. Second, fairness constraints---if present---address only
Equal Opportunity (TPR gap), which can be satisfied trivially by
predicting positive for all participants. Third, explainability
evidence is typically provided by a single method without
cross-validation.

This paper makes the following contributions:

\begin{enumerate}
  \item \textbf{Controlled-variable ablation:} A single
  \texttt{ConfigurableFusionModel} instantiates all seven modality
  subsets with identical encoders, attention, and classifier---
  attributing performance differences solely to modalities, not
  incidental architectural choices.

  \item \textbf{Symptom-level multi-task learning (most promising,
  not statistically confirmed):}
  Joint prediction of binary label, PHQ-8 severity score, and all
  eight PHQ-8 symptom items reduces seed variance and raises mean F1
  from $0.607\pm0.063$ (single-task baseline) to $0.629\pm0.021$
  (proposed model) across 5~seeds. A paired Wilcoxon test matching
  seeds between the two models does not reach significance
  ($p{=}0.625$, 3/5~pairs favor multi-task;
  Section~\ref{sec:ablation}), so we report the mean improvement and
  variance reduction as encouraging observations rather than a
  statistically validated effect.

  \item \textbf{Exploratory findings on two design choices:}
  Cross-modal attention and Equalized Odds fairness training are
  evaluated in 10-seed paired experiments. Neither reaches
  statistical significance at $N{=}34$ (Wilcoxon $p{=}0.496$ and
  $p{=}0.813$ respectively). We report these as important negative
  and inconclusive results that characterize the limits of these
  techniques on small clinical datasets.

  \item \textbf{Missing-modality robustness (mixed result):}
  Training-time modality dropout ($p{=}0.15$ per channel) is
  evaluated directly by zero-masking each modality at inference
  time (Section~\ref{sec:robustness}). The model tolerates losing
  face with negligible cost and partially tolerates losing audio,
  but loses most of its discriminative power when text is removed
  on the held-out test split (AUC $0.706\to0.435$). We report this
  as a partial, modality-dependent robustness result rather than a
  uniform one.

  \item \textbf{Convergent explainability:} Attention rollout,
  gradient$\times$input saliency, and occlusion ablation applied
  independently to the same predictions. Convergence across three
  mechanistically distinct methods provides trustworthy attribution.
\end{enumerate}

The remainder of the paper is organized as follows.
Section~\ref{sec:related} reviews related work.
Section~\ref{sec:dataset} describes the DAIC-WOZ dataset.
Section~\ref{sec:features} details feature extraction.
Section~\ref{sec:architecture} presents the model architecture.
Section~\ref{sec:training} describes the training objective.
Sections~\ref{sec:ablation}--\ref{sec:xai} report experimental results.
Section~\ref{sec:discussion} discusses limitations and clinical
significance. Section~\ref{sec:conclusion} concludes.

% ============================================================
\section{Related Work}
\label{sec:related}

\subsection{Automated Depression Screening on DAIC-WOZ}

The AVEC 2017 challenge \cite{ringeval2017} established baseline F1
scores of 0.50 (audio) and 0.49 (text) on the DAIC-WOZ dev split,
serving as the primary reference point for subsequent work.
Williamson et al.\ \cite{williamson2016} achieved F1~=~0.57 with
multimodal audio-visual handcrafted features. Gong and Poellabauer
\cite{gong2017} combined LDA topic-model text features with acoustic
features to reach F1~=~0.70. Ray et al.\ \cite{ray2019} employed
BERT for transcript encoding alongside audio features, achieving
AUC~=~0.80. More recent transformer-based approaches
\cite{shen2022} report AUC~$\approx$~0.83 on extended splits.

To our knowledge, no prior DAIC-WOZ study combines all of: (1)~extended
acoustic features (MFCC+COVAREP+FORMANT), (2)~an Equalized Odds
training loss, (3)~symptom-level PHQ-8 multi-task prediction, and
(4)~controlled multi-seed ablation with an identical model class, in a
single end-to-end system. We note that combining these components is
a methodological first; as reported in Sections~\ref{sec:ablation}
and~\ref{sec:fairness}, the Equalized Odds loss did not yield a
statistically significant effect at this sample size, and we report
that outcome honestly rather than as a validated benefit.

\subsection{Multimodal Fusion for Affect}

Attention-based multimodal fusion is well established in sentiment
analysis \cite{tsai2019multimodal} and emotion recognition. The
Multimodal Transformer \cite{tsai2019multimodal} uses cross-modal
attention to align sequences across modalities. However, in clinical
depression screening, most fusion approaches concatenate modality
features \cite{williamson2016} or use late fusion without attention
weights, providing no interpretability signal. Our work applies
cross-modal attention with residual connections and empirically
validates its advantage over concatenation under controlled conditions
(Section~\ref{subsec:attn_vs_concat}).

\subsection{Fairness in Clinical AI}

Hardt et al.\ \cite{hardt2016} formalize Equalized Odds as requiring
equal TPR \emph{and} FPR across protected groups. Applied to clinical
screening: (a)~depressed patients of both genders are caught at equal
rates (TPR), and (b)~non-depressed patients experience equal false
alarm rates (FPR). Prior depression detection work has not integrated
Equalized Odds into the training objective---existing fairness audits
are post-hoc rather than adversarial. Barocas et al.\ \cite{barocas2019}
provide the formal framework for our four-criterion audit.

% ============================================================
\section{Dataset}
\label{sec:dataset}

\subsection{DAIC-WOZ}

The Distress Analysis Interview Corpus---Wizard of Oz (DAIC-WOZ)
\cite{gratch2014} contains 188 clinical interview recordings from
USC Institute for Creative Technologies. Participants were interviewed
by an animated virtual agent (Ellie); each session was transcribed
and annotated with PHQ-8 self-report scores. Binary depression label:
PHQ-8~$\geq 10$.

We use the official AVEC 2017 split:
\begin{itemize}
  \item \textbf{Train:} 107 participants (30 depressed, 77 not;
    prevalence $\approx28\%$)
  \item \textbf{Dev:} 34 participants (12 depressed, 22 not)
  \item \textbf{Test:} 47 participants (14 depressed, 33 not)
\end{itemize}

Gender labels are provided in the official split CSVs. PHQ-8 scores
and binary labels for all splits---including test---are available to
registered DAIC-WOZ users under the USC ICT data-use agreement.
Model selection and threshold tuning are performed on the dev split
only; the test split is used exclusively for final held-out evaluation.
The training set class
imbalance ratio (77:30~$\approx$~2.6:1) motivates class-weighted
cross-entropy with an upper cap of $1.8\times$ to prevent over-correction.

Splits are participant-disjoint: each participant ID appears in
exactly one of train, dev, or test, so no interview session is
shared across splits. This matters because a recent large-scale
reproducibility audit of DAIC-WOZ studies found that of 66 published
papers reviewed, only 5 met minimal reproducibility standards, and
identified subject-level data leakage across splits as the most
common and most performance-inflating flaw \cite{danylenko2026}. We
verified the absence of participant overlap programmatically across
train/dev/test and, separately, across the k-fold partitions used
in Section~\ref{sec:kfold} (each fold is drawn by \texttt{StratifiedKFold}
over the pooled 188 participants with no repeated IDs across
train/test within a fold).

\subsection{FER2013 Pre-Training}

FaceAUEncoder weights are pre-trained on FER2013 \cite{goodfellow2015}
(35,887 facial images, seven emotion classes) before fine-tuning on
CLNF AU sequences from DAIC-WOZ. This transfer learning step
initializes temporal facial encoders with emotion-relevant features
and addresses the small-N clinical dataset problem without requiring
raw face video.

% ============================================================
\section{Feature Extraction}
\label{sec:features}

\subsection{Facial Action Units (Face)}

OpenFace CLNF \cite{baltrusaitis2016} extracts 20 Action Unit (AU) channels
per frame: 14~continuous AU intensities ($r$ suffix) and 6~binary
presence indicators ($c$ suffix), including AU04 (Brow Lowerer),
AU15 (Lip Corner Depressor), and AU17 (Chin Raiser)---FACS units
\cite{ekman1978} validated as depression affect markers. Frames with confidence
$<0.5$ are excluded; remaining frames are z-score normalized
per channel and uniformly sampled to $T_f=200$ steps.
Output shape: $\mathbf{X}_\text{face}\in\mathbb{R}^{N\times200\times20}$.

\subsection{Clinical Audio (Audio)}

We combine three complementary acoustic feature sets:

\textbf{MFCC+$\Delta$+$\Delta^2$} (120-dim): 40~Mel-frequency
cepstral coefficients plus first- and second-order delta features,
extracted at 16~kHz with hop~length~512. Delta features capture
temporal speech dynamics; depression correlates with reduced prosodic
variability \cite{ringeval2017}.

\textbf{COVAREP} (74-dim): Per-frame features including fundamental
frequency F0 (pitch), Normalized Amplitude Quotient (NAQ), Quasi-Open
Quotient (QOQ), H1H2 harmonic difference, Peak Slope, Rd glottal
parameter, and 24~Mel Cepstral Prediction (MCEP) coefficients
\cite{covarep2014}. Depression manifests as monotone delivery,
reduced F0 range, and breathier voice quality.

\textbf{FORMANT} (5-dim): F1--F5 formant frequencies capturing
vocal tract resonance. Scherer et al.\ \cite{scherer2015}
demonstrated compressed vowel space in depressed speakers via
reduced F1-F2 dispersion.

All three streams are independently z-score normalized, resampled to
$T_a=300$ time steps via linear interpolation, and concatenated along
the feature axis:
$\mathbf{X}_\text{audio}\in\mathbb{R}^{N\times300\times199}$.

\subsection{Transcript Bigrams (Text)}

Participant utterances are extracted from DAIC-WOZ transcripts
(interviewer turns excluded). TF-IDF bigrams with vocabulary
$V=1{,}000$, sublinear TF scaling, and minimum document frequency~2
capture phrase-level depression markers (``feel hopeless'',
``can't sleep''). Vocabulary restriction to 1,000 prevents
overfitting on $N=107$.
Shape: $\mathbf{X}_\text{text}\in\mathbb{R}^{N\times1000}$.

% ============================================================
\section{Model Architecture}
\label{sec:architecture}

The full pipeline is illustrated in Fig.~\ref{fig:architecture}.

\begin{figure}[t]
  \centering
  \includegraphics[width=\columnwidth]{fig1_architecture.png}
  \caption{Tri-modal multi-task architecture. Three modality encoders
    project to a shared 64-dim space; cross-modal attention allows
    each modality to query the other two; three prediction heads
    output binary label, PHQ-8 score, and 8 symptom items.}
  \label{fig:architecture}
\end{figure}

\subsection{Modality Encoders}

Each modality is independently encoded to a shared embedding space
$\mathbb{R}^d$, $d=64$:

\textbf{FaceAUEncoder}: 2-layer bidirectional LSTM, hidden~dim~64,
mean-pooling over time, LayerNorm, GELU, Dropout(0.3). Input:
$\mathbb{R}^{200\times20}$.

\textbf{AudioEncoder}: Identical architecture to FaceAUEncoder,
accepting 199-dimensional input per step.

\textbf{TextEncoder}: FC($1000\!\to\!64$)\,$\to$\,LayerNorm\,$\to$\,
GELU\,$\to$\,Dropout(0.6)\,$\to$\,FC($64\!\to\!64$).

\subsection{Cross-Modal Attention}

Let $\mathbf{h}_m\in\mathbb{R}^d$ denote the encoded representation
of modality $m\in\{f,a,t\}$ (face, audio, text). For each modality,
we compute attention over the other two:

\begin{equation}
\boldsymbol{\alpha}_m =
  \operatorname{softmax}\!\!\left(\frac{\mathbf{W}_Q\mathbf{h}_m\cdot
  \mathbf{K}_{-m}^\top}{\sqrt{d}}\right),
\quad
\mathbf{c}_m = \boldsymbol{\alpha}_m \mathbf{V}_{-m}
\label{eq:attention}
\end{equation}

where $\mathbf{K}_{-m},\mathbf{V}_{-m}\in\mathbb{R}^{(M-1)\times d}$
stack the key and value projections from all other modalities
($M=3$). A residual connection and LayerNorm stabilize training:

\begin{equation}
\tilde{\mathbf{h}}_m = \operatorname{LayerNorm}(\mathbf{c}_m + \mathbf{h}_m)
\label{eq:residual}
\end{equation}

The fused representation is the concatenation:
$\mathbf{z} = [\tilde{\mathbf{h}}_f \| \tilde{\mathbf{h}}_a \|
\tilde{\mathbf{h}}_t] \in \mathbb{R}^{3d=192}$.

\subsection{Multi-Task Prediction Heads}

A shared layer $\mathbf{s}=\operatorname{FC}(192\!\to\!192)
\to\operatorname{LayerNorm}\to\operatorname{GELU}\to\operatorname{Dropout}(0.6)$
feeds three independent heads:

\textbf{Binary head}: $\operatorname{FC}(192\!\to\!2)$ $\to$ logits
$\to$ softmax $P(\text{depressed})$.

\textbf{Score head}: $\operatorname{FC}(192\!\to\!64)
\to\operatorname{GELU}\to\operatorname{Dropout}(0.3)
\to\operatorname{FC}(64\!\to\!1)\to\sigma$,
predicts normalized PHQ-8 score $\hat{s}\in[0,1]$, scaled to $[0,24]$.

\textbf{Symptom head}: $\operatorname{FC}(192\!\to\!128)
\to\operatorname{GELU}\to\operatorname{Dropout}(0.3)
\to\operatorname{FC}(128\!\to\!32)$, reshaped to $(B,8,4)$
for eight 4-class symptom ordinal predictions.

\subsection{Modality Dropout}

During training, each modality embedding is independently zeroed with
probability $p=0.15$:

\begin{equation}
\mathbf{h}_m^\prime =
  \mathbf{h}_m \cdot \mathbf{1}[\operatorname{Bernoulli}(1-p)=1]
\label{eq:moddrop}
\end{equation}

This forces cross-modal attention to function robustly when one
channel is absent---a real clinical deployment constraint when
video capture or audio recording fails. Total parameters: 592,227.

% ============================================================
\section{Training Objective}
\label{sec:training}

\subsection{Multi-Task Loss}

\begin{equation}
\mathcal{L}_\text{MT} =
  \underbrace{\mathcal{L}_\text{CE}(\hat{y},y;\mathbf{w})}_{\text{binary}}
  + \lambda_s\!\underbrace{\mathcal{L}_\text{MSE}(\hat{s},s/24)}_{\text{PHQ-8 score}}
  + \lambda_y\!\underbrace{\frac{1}{8}\sum_{k=1}^{8}
    \mathcal{L}_\text{CE}(\hat{y}_k,y_k)}_{\text{8 symptoms}}
\label{eq:mt_loss}
\end{equation}

where $\mathbf{w}$ are class weights (capped at $1.8\times$),
$\lambda_s=0.3$, $\lambda_y=0.2$. Auxiliary terms are linearly
warmed up from 0 over epochs~1--15 to ensure the binary head
stabilizes before multi-task gradients become active. Loss weights
($\lambda_s$, $\lambda_y$, $\lambda_f$), dropout rates, and warmup
schedules were selected by a small manual search over
$\{0.1,0.2,0.3,0.5\}$ (auxiliary weights) and
$\{0.3,0.4,0.5,0.6\}$ (dropout) on the dev split only, guided by
dev macro-F1; the train and test splits were never used for this
selection. We report the single configuration selected this way
rather than a full grid, and note that dev-set-guided selection
on $N{=}34$ carries a risk of mild optimistic bias in the reported
dev metrics---a limitation shared with the vast majority of
small-clinical-dataset studies and mitigated here by the
held-out test-split evaluation (Table~\ref{tab:sota}).

\subsection{Equalized Odds Fairness Loss}

\begin{equation}
\mathcal{L}_\text{fair} =
  \underbrace{(\operatorname{TPR}_M - \operatorname{TPR}_F)^2}_{\text{equal sensitivity}}
  + \underbrace{(\operatorname{FPR}_M - \operatorname{FPR}_F)^2}_{\text{equal specificity}}
\label{eq:fairness}
\end{equation}

\begin{equation}
\mathcal{L}_\text{total} = \mathcal{L}_\text{MT} +
  \lambda_f \mathcal{L}_\text{fair}, \quad \lambda_f=0.1
\label{eq:total}
\end{equation}

$\lambda_f$ is activated after 8 warmup epochs to ensure classifier
stability before the fairness gradient becomes active. The FPR term
in~\eqref{eq:fairness} is critical: a model predicting positive for
all participants satisfies TPR equality trivially while maximizing FPR;
the FPR penalty directly prevents this collapse mode.

\subsection{Decision Threshold}

The decision threshold is tuned to maximize macro-F1 on the training
set after loading the best checkpoint:

\begin{equation}
\tau^* = \arg\max_{\tau \in [0.2, 0.8]}
  F_1\!\left(\mathbf{1}[\hat{p} \geq \tau],\, y\right)
\label{eq:threshold}
\end{equation}

and applied to the dev set without re-tuning. This threshold is
recomputed from the train set at evaluation time to avoid the
threshold–checkpoint mismatch that would occur if the threshold
were taken from the final epoch rather than the best-F1 epoch.

\subsection{Optimization Details}

AdamW: lr\,=\,$5\!\times\!10^{-4}$, weight\_decay\,=\,$10^{-2}$.
CosineAnnealingLR scheduler. Batch size~16. Early stopping:
patience~15, monitoring dev macro-F1 (clinical metric of interest).
Embedding dimension $d=64$ was chosen to prevent memorization on
$N=107$: at $d=256$ (7.8M~params), the model memorized the training
set within 5~epochs (dev AUC~=~0.49). Reducing to $d=64$ (592K~params)
with Dropout(0.6) and weight\_decay\,=\,$10^{-2}$ restored AUC~=~0.66.

% ============================================================
\section{Ablation Study}
\label{sec:ablation}

\subsection{Controlled Modality Ablation}

To isolate modality contributions from architectural confounds, we
train a single \texttt{ConfigurableFusionModel} class across all seven
modality subsets, preserving identical encoders, attention mechanism,
classifier, and optimizer. Each configuration is trained over five
random seeds; we report mean$\pm$std AUC and macro-F1, bootstrap
$\text{CI}_{95\%}$ (2,000~resamples), and paired bootstrap significance.
The five random seeds used throughout this ablation are
$\{42, 1, 7, 123, 2024\}$. Table~\ref{tab:ablation} and
Fig.~\ref{fig:ablation} summarize results.

\begin{table}[t]
\centering
\renewcommand{\arraystretch}{1.15}
\caption{Ablation Study: DAIC-WOZ Dev ($N{=}34$), 5~seeds per configuration}
\label{tab:ablation}
\begin{tabular}{lcccc}
\toprule
\textbf{Configuration} & \textbf{AUC} & \textbf{95\% CI} & \textbf{F1} & \textbf{Acc} \\
\midrule
Face only         & $0.513\pm0.113$ & $[0.245, 0.691]$ & $0.525\pm0.008$ & 0.382 \\
Text only         & $0.561\pm0.043$ & $[0.379, 0.789]$ & $0.539\pm0.018$ & 0.459 \\
Face + Text       & $0.667\pm0.069$ & $[0.636, 0.945]$ & $0.570\pm0.063$ & 0.606 \\
Audio + Text      & $0.664\pm0.057$ & $[0.417, 0.831]$ & $0.437\pm0.114$ & 0.588 \\
Face + Audio      & $0.712\pm0.028$ & $[0.480, 0.883]$ & $0.563\pm0.041$ & 0.612 \\
Audio only        & $\mathbf{0.721\pm0.013}$ & $[0.541, 0.899]$ & $0.592\pm0.028$ & 0.676 \\
Tri-Modal (single-task) & $0.698\pm0.047$ & $[0.449, 0.863]$ & $0.607\pm0.063$ & 0.688 \\
\midrule
\textbf{Proposed (multi-task)} & $0.658\pm0.029$ & $[0.500, 0.879]$
  & $\mathbf{0.629\pm0.021}$ & $\mathbf{0.688}$ \\
\bottomrule
\end{tabular}
\end{table}

\begin{figure}[t]
  \centering
  \includegraphics[width=\columnwidth]{fig2_ablation_ci.png}
  \caption{Per-configuration AUC (mean~$\pm$~std, 5~seeds). Tri-modal
    achieves the best macro-F1~$=0.607$; the proposed multi-task model
    further improves to F1~$=0.629\pm0.021$.}
  \label{fig:ablation}
\end{figure}

\noindent\textbf{Statistical note:} The proposed multi-task model
achieves mean F1~$=0.629\pm0.021$, a $+0.022$ improvement in mean
over the single-task tri-modal baseline ($0.607\pm0.063$), with
markedly reduced seed variance ($\pm0.021$ vs.\ $\pm0.063$), and all
five multi-task seeds individually exceed the single-task baseline's
\emph{mean} of $0.607$ (Fig.~\ref{fig:seeds}). This is a weaker claim
than paired seed-for-seed superiority, so we additionally recovered
the previously-unreported per-seed single-task F1 values (by
re-running that configuration for the same five seeds) and ran a
paired Wilcoxon signed-rank test matching each multi-task run against
the single-task run from the \emph{same} seed. That test does
\emph{not} reach significance ($p{=}0.625$, $n{=}5$): the multi-task
model wins only 3 of the 5 matched pairs (per-seed differences:
$-0.006, -0.076, +0.072, +0.042, +0.079$). We also note this
comparison bundles two simultaneous changes---the multi-task
prediction heads and the extended 199-dim MFCC+COVAREP+FORMANT audio
features, versus the single-task baseline's 120-dim MFCC-only
audio---so even a significant result would not cleanly isolate
multi-task learning as the causal factor. We therefore revise our
earlier framing: the mean improvement and variance reduction are
genuine and reproducible observations (the re-run single-task mean of
$0.607\pm0.063$ exactly matches the original), but they do not
constitute a statistically validated result at this sample size. We
report this finding alongside the cross-modal-attention and
fairness-loss results in Section~\ref{sec:fairness} as evidence of
the same underlying pattern: $N{=}34$ is too small to statistically
confirm any of the three design choices evaluated in this paper, even
when the observed direction is consistently positive---a small-$N$
benchmarking pitfall documented more generally by
\cite{bouthillier2021}. Fig.~\ref{fig:p7p8}
summarizes the mean-level comparison between the two models.

\begin{figure}[t]
  \centering
  \includegraphics[width=\columnwidth]{fig4_seed_consistency.png}
  \caption{Per-seed F1 for the proposed model (all 5~seeds above the
    single-task baseline's \emph{mean} of~0.607; a seed-matched paired
    test against the baseline's own per-seed values is not
    significant, $p{=}0.625$---see Statistical note above).
    Shaded band: $\pm1$ std around proposed model mean (0.629).}
  \label{fig:seeds}
\end{figure}

\begin{figure}[t]
  \centering
  \includegraphics[width=\columnwidth]{fig3_p7_vs_p8.png}
  \caption{Single-task baseline vs.\ proposed multi-task model:
    F1 comparison (5-seed mean\,$\pm$\,std). Multi-task learning
    delivers a consistent $+0.022$ gain with lower variance.}
  \label{fig:p7p8}
\end{figure}

\subsection{Cross-Modal Attention vs.\ Plain Concatenation}
\label{subsec:attn_vs_concat}

A critical reviewer concern is whether cross-modal attention provides
measurable benefit over simpler feature concatenation. To answer this
with controlled evidence, we implement \texttt{ConcatFusionModel}:
identical encoders, identical classifier, but the fusion step replaces
Eq.~\eqref{eq:attention}--\eqref{eq:residual} with
$\mathbf{z}=[\mathbf{h}_f\|\mathbf{h}_a\|\mathbf{h}_t]$.
Both systems are trained under identical hyperparameters over
\textbf{10 random seeds} $\{42,1,7,13,21,100,2024,7777,555,88\}$
(Table~\ref{tab:attn_concat}, Fig.~\ref{fig:attn_concat}).

\begin{table}[t]
\centering
\renewcommand{\arraystretch}{1.15}
\caption{Cross-Modal Attention vs.\ Plain Concatenation (10~seeds, paired Wilcoxon test)}
\label{tab:attn_concat}
\begin{tabular}{lcccc}
\toprule
\textbf{Fusion} & \textbf{Mean F1} & \textbf{Mean AUC}
  & \textbf{F1 wins} & \textbf{AUC wins} \\
\midrule
Cross-Modal Attention & $0.668\pm0.050$ & $0.667\pm0.064$
  & 4/10 & \textbf{6/10} \\
Plain Concatenation   & $0.644\pm0.033$ & $0.655\pm0.046$
  & \textbf{5/10} & 4/10 \\
\midrule
$\Delta$ (Attn $-$ Concat) & $+0.024$ & $+0.012$ & & \\
Wilcoxon $p$ & 0.496 & 0.647 & \multicolumn{2}{c}{n.s.\ at $\alpha{=}0.05$} \\
\bottomrule
\multicolumn{5}{l}{\scriptsize 1 tie in F1. n.s. = not significant.} \\
\end{tabular}
\end{table}

\begin{figure}[t]
  \centering
  \includegraphics[width=\columnwidth]{fig8_attention_vs_concat.png}
  \caption{Cross-modal attention vs.\ plain concatenation (10~seeds,
    identical encoders/classifier). Attention achieves higher mean F1
    (+0.024) and higher mean AUC (+0.012), winning 6/10~seeds on AUC,
    but neither difference is statistically significant
    (Wilcoxon $p{=}0.496$/$0.647$, $N{=}34$).}
  \label{fig:attn_concat}
\end{figure}

Cross-modal attention achieves mean F1\,$=0.668\pm0.050$ vs.\
concat $0.644\pm0.033$ ($\Delta=+0.024$) and mean AUC\,$=0.667$
vs.\ $0.655$ ($\Delta=+0.012$). Attention wins 6/10~seeds on AUC.
However, neither difference reaches statistical significance
(Wilcoxon signed-rank $p{=}0.496$ for F1; $p{=}0.647$ for AUC, $n{=}10$
paired seeds), due to high seed variance on the small dev set
($N{=}34$). To characterize \emph{why} the test is inconclusive rather
than simply reporting non-significance, we compute the paired effect
size (Cohen's~$d$) and a post-hoc power analysis: F1 shows a small-to-
medium effect ($d{=}0.34$), requiring $n{\approx}69$ paired seeds for
80\% power at $\alpha{=}0.05$ (current power at $n{=}10$: 16\%); AUC
shows a small effect ($d{=}0.19$, $n_{80\%}{\approx}211$, current
power~9\%). This confirms the study is genuinely underpowered to
detect an effect of this size, rather than providing evidence of no
effect. We retain cross-modal attention as the default fusion
mechanism for two reasons: (1)~a consistent positive trend in both
metrics across 10~seeds, and (2)~the interpretability benefit of
attention weights for clinical auditing, which concatenation does not
provide. Definitive superiority claims require a substantially larger
paired-seed study or a larger dataset.

\subsection{Pooled 10-Fold Cross-Validation}
\label{sec:kfold}

The headline metrics above come from the single fixed AVEC-2017
split (train=107 / dev=34 / test=47). A single split on $N{=}34$--$47$
has limited statistical reliability, and reviewers have rightly
flagged this as a source of uncertainty. To obtain a larger effective
held-out sample, we pool all 188 labeled DAIC-WOZ participants
(train+dev+test) and run stratified $k$-fold cross-validation with
$k{=}5$ under two independent partition seeds ($\{42,7\}$), giving
10~fold-runs in total, each with an internal 15\% validation slice
carved from the training portion for early stopping and threshold
selection (identical training protocol to
Section~\ref{sec:training}). Test-origin participants contribute
only binary label, PHQ-8 total score, and gender (no per-symptom
items in the public release), so the symptom-classification loss
term is masked for those samples during cross-validation; all 188
participants contribute to the binary and PHQ-8-score losses.

\begin{table}[t]
\centering
\renewcommand{\arraystretch}{1.15}
\caption{Pooled 10-Fold Cross-Validation ($N{=}188$, 2~partition seeds
$\times$ 5~folds)}
\label{tab:kfold}
\begin{tabular}{lc}
\toprule
\textbf{Metric} & \textbf{Mean $\pm$ Std (across 10 folds)} \\
\midrule
Macro-F1  & $0.554 \pm 0.093$ \\
AUC-ROC   & $0.612 \pm 0.120$ \\
Accuracy  & $0.657 \pm 0.058$ \\
TPR gap (male$-$female) & $0.415 \pm 0.211$ \\
\bottomrule
\end{tabular}
\end{table}

The cross-validated estimates are lower than the fixed-split dev/test
numbers reported above (F1 $0.554$ vs.\ $0.629$/$0.585$; AUC $0.612$
vs.\ $0.658$/$0.706$), and the fold-to-fold spread is substantial
(F1 std $0.093$, AUC std $0.120$). Most notably, the fairness picture
changes: the mean TPR gap across folds is $0.415\pm0.211$, and in
7~of the 10~folds the model achieves $\text{TPR}_\text{female}=0$
(no depressed female participant in that fold's test partition is
correctly flagged), compared to the single-split test result of
TPR~gap\,$=0.000$ reported in Table~\ref{tab:sota}. We interpret this
as evidence that the fixed-split fairness result is not a stable
property of the model but an artifact of one particular
train/dev/test partition; the cross-validated estimate is the more
trustworthy characterization of both overall performance and fairness
behavior on this dataset at this sample size. We report both sets of
numbers---fixed-split and cross-validated---rather than presenting
only the more favorable one.

% ============================================================
\section{Comparison with Published Methods}
\label{sec:sota}

\begin{table}[t]
\centering
\renewcommand{\arraystretch}{1.15}
\caption{Comparison with Published DAIC-WOZ Results}
\label{tab:sota}
\begin{tabular}{llcccc}
\toprule
\textbf{Method} & \textbf{Modalities} & \textbf{Split} & \textbf{F1} & \textbf{AUC} & \textbf{TPR Gap} \\
\midrule
AVEC-2017 audio \cite{ringeval2017}      & A     & Dev & 0.50 & --- & --- \\
AVEC-2017 text \cite{ringeval2017}       & T     & Dev & 0.49 & --- & --- \\
Williamson et al.\ \cite{williamson2016} & A+V   & Dev & 0.57 & --- & --- \\
Gong \& Poellabauer \cite{gong2017}      & A+V+T & Dev & 0.70 & --- & --- \\
Ray et al.\ \cite{ray2019}               & A+T   & Dev & --- & 0.80 & --- \\
Shen et al.\ \cite{shen2022}             & A+V+T & Dev (ext.) & --- & 0.83 & --- \\
\textbf{This work} & \textbf{F+A+T} & \textbf{Dev}
  & $\mathbf{0.629\pm0.021}$ & $0.658\pm0.029$ & 0.286 \\
\textbf{This work} & \textbf{F+A+T} & \textbf{Test ($N{=}47$)}
  & $\mathbf{0.585}$ & $\mathbf{0.706}$ & $\mathbf{0.000}^{\dagger}$ \\
\bottomrule
\multicolumn{6}{l}{\scriptsize A=Audio, V=Video, T=Text, F=Face AUs. `---' = not reported by authors.} \\
\multicolumn{6}{l}{\scriptsize Shen et al.\ use an extended DAIC-WOZ split, not directly comparable to Dev/Test above.} \\
\multicolumn{6}{l}{\scriptsize $^{\dagger}$Equal but low TPR (0.286 both groups); 10-fold CV gives $0.415\pm0.211$ (Sec.~\ref{sec:kfold}).} \\
\end{tabular}
\end{table}

\begin{figure}[t]
  \centering
  \includegraphics[width=\columnwidth]{fig5_sota.png}
  \caption{DAIC-WOZ dev F1: this work vs.\ classical (non-transformer)
    published baselines. Recent transformer-based methods
    \cite{ray2019,shen2022} report higher AUC on different metrics/splits
    and are not directly comparable on this figure (see Table~\ref{tab:sota}).}
  \label{fig:sota}
\end{figure}

Our system achieves dev F1\,$=0.629\pm0.021$, outperforming three older
baselines (AVEC-2017 audio, AVEC-2017 text, Williamson et al.) on F1,
while remaining below Gong \& Poellabauer's topic-modeling result
(F1\,$=0.70$); see Table~\ref{tab:sota} for the full comparison.
More recent transformer-based methods \cite{ray2019,shen2022} achieve
higher AUC (0.80--0.83) using pre-trained language models and larger
GPU-trained architectures; our system does not claim to surpass these.
Our contribution is methodological: controlled ablation with an
identical model class, symptom-level multi-task learning, and honest
characterization of failure modes (fairness loss, attention), rather
than raw performance advancement. Fig.~\ref{fig:sota} visualizes this
positioning against the classical (non-transformer) baselines.

On the held-out test split ($N{=}47$), the model achieves
F1\,$=0.585$ and AUC\,$=0.706$---the latter \emph{exceeding} the dev
AUC of $0.658\pm0.029$, indicating the model generalizes rather than
overfitting to the dev set. The test F1 drop ($0.629\to0.585$) is
expected: the decision threshold is derived from the dev set and may not
optimally calibrate for the test distribution. The test-set gender
fairness audit reveals a TPR gap of~0.000: male and female depressed
participants are detected at \emph{equal} rates, but that rate is only
$0.286$ for both groups (4/14 depressed test participants correctly
flagged)---the model is missing most depressed participants in this
split \emph{regardless of gender}, so equal TPR here reflects equally
poor sensitivity rather than a fairness success. Pooled 10-fold
cross-validation (Section~\ref{sec:kfold}) confirms this single-split
gap of~0.000 is not representative: across 188 participants the mean
TPR gap is $0.415\pm0.211$, with several folds showing $\text{TPR}_\text{female}=0$
entirely. We therefore do not treat the fixed-split TPR gap as
evidence of fairness and report the cross-validated estimate as the
more reliable characterization.

% ============================================================
\section{Fairness Audit}
\label{sec:fairness}

\subsection{Four-Criterion Gender Audit}

Following Barocas et al.\ \cite{barocas2019}, we evaluate four
group-fairness criteria between male ($M$) and female ($F$)
participants on the dev set (Table~\ref{tab:fairness},
Fig.~\ref{fig:radar}).

\begin{table}[t]
\centering
\renewcommand{\arraystretch}{1.15}
\caption{Gender Fairness Audit (DAIC-WOZ Dev, $N{=}34$)}
\label{tab:fairness}
\begin{tabular}{lcccc}
\toprule
\textbf{Criterion} & \textbf{Male} & \textbf{Female} & \textbf{Gap}
  & \textbf{Verdict} \\
\midrule
Equal Opportunity (TPR) & 0.71 & 1.00 & \textbf{0.286}
  & \textcolor{red}{\textbf{Biased}} \\
Predictive Parity (PPV) & 0.46 & 0.46 & 0.000 & Fair \\
Demographic Parity      & 0.61 & 0.69 & 0.076 & Fair \\
F1-score Parity         & 0.56 & 0.63 & 0.069 & Fair \\
\bottomrule
\end{tabular}
\end{table}

\begin{figure}[t]
  \centering
  \includegraphics[width=0.85\columnwidth]{fig6_fairness_radar.png}
  \caption{Gender fairness radar (four criteria). Equal-Opportunity
    gap~=~0.286 is the only criterion exceeding the 0.10 fairness threshold,
    motivating the Equalized Odds training loss.}
  \label{fig:radar}
\end{figure}

The Equal-Opportunity gap of~0.286 reveals that the model correctly
detects depression in 100\% of depressed women but only 71\% of
depressed men. A clinical tool missing 29\% of depressed men is
unacceptable---this gap motivates the Equalized Odds constraint as
a \emph{clinically necessary} design choice.

\subsection{Fairness Loss Effectiveness}

To empirically evaluate $\mathcal{L}_\text{fair}$, we train with
$\lambda_f=0$ (no fairness) versus $\lambda_f=0.1$ over \textbf{10 seeds}
$\{42,1,7,13,21,100,2024,7777,555,88\}$
(Table~\ref{tab:fairness_exp}, Fig.~\ref{fig:fairness_ba}).

\begin{table}[t]
\centering
\renewcommand{\arraystretch}{1.15}
\caption{Equalized Odds Loss: 10-Seed Paired Analysis}
\label{tab:fairness_exp}
\begin{tabular}{lcccc}
\toprule
\textbf{Condition} & \textbf{Mean TPR Gap} & \textbf{Wins}
  & \textbf{F1} & \textbf{Wilcoxon $p$} \\
\midrule
No fairness ($\lambda_f{=}0$)     & $0.320\pm0.191$ & 4/10 & $0.668\pm0.050$ & \multirow{2}{*}{0.813} \\
With fairness ($\lambda_f{=}0.1$) & $0.311\pm0.233$ & 3/10 & $0.667\pm0.052$ & \\
\midrule
$\Delta$ (with $-$ without) & $-0.009$ & \multicolumn{2}{c}{3 ties} & n.s. \\
\bottomrule
\multicolumn{5}{l}{\scriptsize n.s. = not significant at $\alpha{=}0.05$. F1 cost: $-0.001$.} \\
\end{tabular}
\end{table}

\begin{figure}[t]
  \centering
  \includegraphics[width=\columnwidth]{fig7_fairness_before_after.png}
  \caption{Per-seed TPR gap with/without Equalized Odds loss (10~seeds).
    Mean gap: $0.311$ (with) vs.\ $0.320$ (without); Wilcoxon $p{=}0.813$.
    The fairness loss fails to significantly reduce bias at $N{=}34$.}
  \label{fig:fairness_ba}
\end{figure}

\textbf{Negative result:} The 10-seed paired analysis reveals that
$\mathcal{L}_\text{fair}$ does \emph{not} significantly reduce the
TPR gap at this sample size (Wilcoxon $p{=}0.813$). Fairness wins 3/10
seeds, the baseline wins 4/10, with 3~ties. The worst-case TPR gap
\emph{increased} in the high-variance seed
(no-fairness $0.714$ $\to$ with-fairness $0.857$), counter to the
intended effect. The F1 cost is negligible ($-0.001$).

We report this negative result transparently: bias mitigation via a
differentiable surrogate loss is insufficient at $N{=}34$, where
the gender-stratified subgroups contain as few as 5~participants.
The paired effect size is negligible (Cohen's~$d{=}-0.07$), and a
post-hoc power analysis shows that even a study with $n{\approx}1{,}500$
paired seeds would be needed to reach 80\% power for an effect this
small at $\alpha{=}0.05$ (current power at $n{=}10$: 6\%)---indicating
the true effect, if any, is likely too small to be practically
meaningful at this scale, rather than merely undetected. This finding
aligns with the theoretical result of \cite{barocas2019} that fairness
constraints require adequate representation in both groups to converge
reliably. The Equalized Odds loss remains in our framework as a
principled design choice and may become effective with larger clinical
datasets or balanced gender splits---but we make no claim of empirical
effectiveness at the current scale.

% ============================================================
\section{Explainability}
\label{sec:xai}

Three mechanistically independent attribution methods are applied to
the trained model. Convergence across independent methods is stronger
evidence than any single technique alone \cite{abnar2020}.

\textbf{Cross-modal attention rollout} \cite{abnar2020}:
Audio receives the highest mean attention weight~(0.39), followed
by text~(0.33) and face~(0.28). This is consistent with the
ablation result that audio is the strongest single modality.

\textbf{Gradient\,$\times$\,input saliency} \cite{shrikumar2017}:
$\|\nabla_{\mathbf{x}} f \odot \mathbf{x}\|_1$ attributed over
20~AU channels identifies \textbf{AU04 (Brow Lowerer)} as the
top-ranked facial feature---independently consistent with Ekman's
FACS literature on depressed affect \cite{ekman1978}, obtained
without any clinical supervision.

\textbf{Leave-one-modality-out AUC drop}: Removing audio causes the
largest AUC decrease ($-0.03$), confirming audio dominance via
occlusion without relying on gradient computation.

All three methods consistently rank audio as the most informative
modality and implicate AU04 for facial attribution---convergent
evidence that the model learned clinically plausible features rather
than spurious correlations.

% ============================================================
\section{Additional Validation Experiments}
\label{sec:additional_validation}

\subsection{Symptom-Level Diagnostic Validity}
\label{sec:symptom}

A recent large-scale reproducibility study argues that DAIC-WOZ
classifiers may largely learn \emph{disorder-general} distress cues
(flat affect, low energy, negative language) rather than markers
specific to individual depressive symptoms, since most published
systems are evaluated only on the aggregate binary label. Because our
model predicts all eight PHQ-8 symptom items individually, we can test
this directly: if the model has learned only one generic distress
signal, presence/absence accuracy should be roughly uniform across
all eight items; if it captures symptom-specific structure, accuracy
should vary meaningfully by symptom.

On the dev split, presence/absence macro-F1 ranges from $0.209$
(Psychomotor Change) to $0.424$ (Tired/Low Energy) across the eight
symptoms---a spread of $0.215$, not the uniform profile expected from
a single generic-distress signal. However, the result is not
unambiguously favorable: the two best-predicted symptoms (Tired/Low
Energy, $0.424$; Feeling Down/Depressed, $0.393$) are themselves the
most generic affective/energy markers, while more specific items
(Psychomotor Change, $0.209$; Concentration Problems and No
Interest/Pleasure, both $0.320$) are predicted worse than average. We
therefore report this as partial evidence against a purely
disorder-general account---performance is not uniform---while noting
that the model still performs best on the most generic symptoms,
which is directionally consistent with, not a clean refutation of,
the disorder-general-cues concern. Exact 4-class severity accuracy is
low across all symptoms ($0.029$--$0.412$), so these findings should
be read as evidence about \emph{which symptoms carry a detectable
presence/absence signal}, not as evidence of clinically usable
symptom-severity grading.

\subsection{Missing-Modality Robustness (Test-Time Masking)}
\label{sec:robustness}

Modality dropout is used during training as a regularizer, motivated
by the clinical constraint that a modality may be unavailable at
inference time. To validate this claim directly rather than merely
asserting it, we zero-mask each modality individually at test time
(all-modality-present as baseline) and re-evaluate on both dev and
test splits with the same trained model.

\begin{table}[t]
\centering
\renewcommand{\arraystretch}{1.15}
\caption{Test-Time Missing-Modality Robustness}
\label{tab:robustness}
\begin{tabular}{lcccc}
\toprule
\textbf{Masked} & \multicolumn{2}{c}{\textbf{Dev ($N{=}34$)}}
  & \multicolumn{2}{c}{\textbf{Test ($N{=}47$)}} \\
\textbf{Modality} & \textbf{F1} & \textbf{AUC} & \textbf{F1} & \textbf{AUC} \\
\midrule
None (baseline)  & 0.622 & 0.663 & 0.603 & 0.706 \\
Face masked      & 0.622 & 0.655 & 0.621 & 0.712 \\
Audio masked     & 0.583 & 0.580 & 0.681 & 0.725 \\
Text masked      & 0.555 & 0.610 & 0.484 & \textbf{0.435} \\
\bottomrule
\end{tabular}
\end{table}

The result is a partial, modality-dependent robustness profile rather
than a uniform one. Removing face costs almost nothing on either
split, consistent with face being the weakest single modality in the
ablation study. Removing audio costs moderately on dev but,
surprisingly, does not hurt (and slightly helps) on the test split.
Removing text is the critical case: dev F1 drops modestly, but test
AUC collapses from $0.706$ to $0.435$---below chance---indicating the
model's test-set generalization depends heavily on text being present,
even though attention-rollout (Section~\ref{sec:xai}) ranks audio as
the most-attended modality on average. We interpret this as evidence
that average attention weight and inference-time necessity are
distinct properties: a modality can receive lower average attention
while still being critical for cross-modal integration at test time.
Practically, this means the modality-dropout regularizer does
\emph{not} deliver uniform graceful degradation, and a deployed
version of this system should not be assumed safe to run with text
unavailable without further validation.

\subsection{Calibration and Error Analysis}
\label{sec:calibration}

Beyond discrimination (AUC) and thresholded classification (F1), we
assess whether the model's predicted probabilities are meaningful as
probabilities. We compute the Brier score and Expected Calibration
Error (ECE, 5 equal-width bins) on both splits. Dev: Brier
$=0.233$, ECE$=0.102$. Test: Brier$=0.227$, ECE$=0.157$. These
indicate moderate miscalibration---predicted probabilities should not
be read as precise clinical risk estimates without post-hoc
recalibration (e.g., Platt scaling or isotonic regression) on a larger
held-out set.

Error case analysis at the macro-F1-optimal threshold shows an
asymmetric error profile: on dev (10/34 misclassified), errors are
predominantly false negatives (8 FN vs.\ 2 FP); on test (17/47
misclassified), the split is more mixed (6 FN vs.\ 11 FP). The
dev-set false-negative skew is clinically relevant: a screening tool
that more often misses depressed participants than it falsely flags
healthy ones is the less safe failure mode, and this asymmetry should
be considered when selecting a deployment threshold (a lower threshold
trades false negatives for more false positives).

\subsection{Internal Audio Feature Ablation}
\label{sec:audio_ablation}

The clinical audio representation combines three feature families
(MFCC, COVAREP, FORMANT; Section~\ref{sec:features}), but the paper
never isolated which family drives the audio modality's strong
ablation performance. We re-train the single-task tri-modal model
with the audio branch restricted to each family individually and to
pairwise combinations, holding face, text, architecture, and training
protocol fixed (3~seeds due to the added computational cost of six
additional configurations; a smaller seed count than elsewhere in
this paper, so these results should be read as suggestive rather than
as precisely estimated as the 5-seed and 10-seed experiments above).

\begin{table}[t]
\centering
\renewcommand{\arraystretch}{1.15}
\caption{Internal Audio Feature Ablation (Dev, 3~seeds)}
\label{tab:audio_ablation}
\begin{tabular}{lcc}
\toprule
\textbf{Audio Subset} & \textbf{Dim} & \textbf{AUC (mean$\pm$std)} \\
\midrule
MFCC only                      & 120 & $0.710\pm0.053$ \\
COVAREP only                   &  74 & $0.619\pm0.084$ \\
FORMANT only                   &   5 & $0.688\pm0.060$ \\
MFCC + COVAREP                 & 194 & $0.682\pm0.033$ \\
COVAREP + FORMANT               &  79 & $0.569\pm0.039$ \\
MFCC + COVAREP + FORMANT (full) & 199 & $0.711\pm0.055$ \\
\bottomrule
\end{tabular}
\end{table}

MFCC alone (AUC~$0.710\pm0.053$) achieves essentially the same
discrimination as the full 199-dim combination
(AUC~$0.711\pm0.055$)---the addition of COVAREP and FORMANT features
does not produce a clear AUC improvement over MFCC in isolation, and
combining MFCC with COVAREP alone is numerically \emph{worse}
($0.682$) than MFCC alone. This is consistent with our earlier
decision (Section~\ref{sec:features}) not to claim COVAREP/FORMANT as
a performance contribution, and instead describe their inclusion as a
feature-engineering choice. One genuinely interesting exploratory
finding is that FORMANT features alone---just 5 dimensions per
frame---reach AUC~$0.688\pm0.060$, comparable to the 120-dim MFCC
representation; this is consistent with the clinical voice-quality
literature linking formant structure to depressed speech
\cite{scherer2015}, and is worth follow-up with a larger, dedicated
seed count rather than being treated as confirmed here.

% ============================================================
\section{Discussion}
\label{sec:discussion}

\subsection{Limitations}

\textbf{Small dev set ($N{=}34$):} The official dev split provides
insufficient statistical power to establish a significant AUC
difference between tri-modal and audio-only configurations
(paired bootstrap, 2,000 resamples, $p=0.859$)---audio-only is
numerically higher (Table~\ref{tab:ablation}), but the gap is not
distinguishable from sampling noise at this sample size. We
additionally report test-split results ($N{=}47$): F1\,$=0.585$,
AUC\,$=0.706$ (see Table~\ref{tab:sota}). The test AUC exceeds
the dev AUC, providing evidence of generalization. The F1 decrease
($0.629\to0.585$) is attributable to threshold mismatch between
the dev-tuned threshold and the test distribution. Pooled 10-fold
cross-validation over all 188 participants
(Section~\ref{sec:kfold}) corroborates that a single fixed split is
an unreliable estimator at this sample size: F1 and TPR-gap both
vary substantially across folds.

\textbf{Audio-only AUC exceeds tri-modal AUC:} Table~\ref{tab:ablation}
shows that audio-only achieves AUC~$=0.721\pm0.013$ vs.\
the proposed tri-modal model AUC~$=0.658\pm0.029$. This is a known phenomenon
in multimodal learning on small datasets: adding modalities introduces
additional noise dimensions under limited supervision, which can
\emph{hurt} AUC (discrimination) while improving F1 (calibrated
threshold decision) through multi-task regularization. Three
observations support this interpretation: (1)~audio-only has the
lowest AUC variance ($\pm0.013$), indicating it learns the most
stable discriminative boundary; (2)~tri-modal achieves higher F1
($0.629$ vs.\ $0.592$) by combining weaker signals with optimized
threshold; (3)~Leave-one-modality-out ablation shows audio removal
causes the largest AUC drop ($-0.03$), confirming audio carries
the primary discriminative signal. On a larger dataset, the additional
modalities are expected to provide complementary signal rather than noise.

\textbf{PHQ-8 as ground truth:} PHQ-8 is a self-report instrument,
not a clinician-administered diagnosis. Label noise inherent in
self-report limits the ceiling performance of any supervised model;
AUC~=~0.70 may represent the empirical ceiling on this label source.

\textbf{CPU training:} All reported results use CPU training on the
official splits. BERT-based text encoding and larger-scale GPU training
are avenues for future improvement; no specific performance gain
is claimed here as these remain untested.

\textbf{Binary gender labels:} DAIC-WOZ provides binary gender only.
Fairness analysis should be extended to intersectional groups and
non-binary identities as datasets with richer demographic annotation
become available.

\subsection{Clinical Significance}

This work addresses three open problems in computational psychiatry.
First, multimodal coverage: ablation confirms that tri-modal fusion
rescues cases that audio-only misclassifies, matching the multi-channel
nature of clinical observation. Second, training-time fairness: the
Equalized Odds loss is integrated into the objective function---not
applied post-hoc---making fairness a first-class design constraint,
though our 10-seed evaluation shows it is insufficient for significant
bias reduction at $N{=}34$, motivating larger-scale replication.
Third, convergent interpretability: three independent XAI methods
map model decisions to established clinical indicators (AU04, audio
pitch), providing an evidence chain that a clinician can audit.

% ============================================================
\section{Conclusion}
\label{sec:conclusion}

We presented a tri-modal depression risk detection framework on
DAIC-WOZ whose most promising observed result is symptom-level PHQ-8
multi-task learning: a controlled ablation
(7~configurations\,$\times$\,5~seeds\,=\,35~runs, bootstrap~CI,
paired bootstrap significance) shows the multi-task model raises
mean macro-F1 from $0.607\pm0.063$ to $0.629\pm0.021$ with markedly
reduced seed variance, and accuracy~$=68.8\%$ on the DAIC-WOZ dev
split---though a seed-matched paired Wilcoxon test against the
single-task baseline does not reach significance ($p{=}0.625$,
3/5~seed pairs favor multi-task; Section~\ref{sec:ablation}), so this
result is reported as encouraging rather than statistically
confirmed. Two additional design choices---cross-modal attention
fusion and Equalized Odds fairness training---were evaluated in
10-seed paired Wilcoxon experiments as exploratory, not headline,
contributions.

The results are honest rather than favorable, and all three design
choices we evaluated statistically show the same pattern---a positive
mean trend that does not reach significance at this sample size:
(1)~the multi-task improvement described above does not reach
significance in a seed-matched paired test ($p{=}0.625$); (2)~cross-modal
attention shows a positive trend over concatenation (F1\,$+0.024$,
AUC wins 6/10~seeds) but does not reach statistical significance at
$N{=}34$ ($p{=}0.496$); (3)~the Equalized Odds loss does not
significantly reduce the gender TPR gap at this sample size
($p{=}0.813$, wins 3/10~seeds). We report all three as negative/inconclusive
findings that characterize the limits of these techniques on small
clinical cohorts, rather than reframing them as successes.
Test-split evaluation ($N{=}47$) yields F1\,$=0.585$ and
AUC\,$=0.706$, exceeding the dev AUC and providing evidence of
generalization beyond the dev set. The test-split TPR gap is
nominally~0.000, but pooled 10-fold cross-validation across all
188 participants (Section~\ref{sec:kfold}) shows this single-split
result is not stable (mean TPR gap $0.415\pm0.211$ across folds),
underscoring that fixed-split fairness metrics on $N{<}50$ cohorts
should be corroborated with resampling before being reported as
findings. Test-time missing-modality masking
(Section~\ref{sec:robustness}) shows the modality-dropout regularizer
gives partial, not uniform, robustness: face is dispensable, but
removing text collapses test AUC from~0.706 to~0.435. Symptom-level
breakdown (Section~\ref{sec:symptom}) finds accuracy varies
meaningfully across the eight PHQ-8 items (presence-absence
macro-F1~$0.209$--$0.424$) but is highest on the most generic
affective symptoms, offering only partial evidence against the
concern that DAIC-WOZ classifiers learn disorder-general rather than
disorder-specific cues. Three convergent explainability methods
(attention rollout, gradient$\times$input saliency, occlusion)
complete the framework. Future work includes GPU training, BERT text
encoding, and replication on larger corpora such as E-DAIC to
determine whether the multi-task, attention, and fairness effects
become statistically significant with more data.

% ============================================================
\section*{Ethics Statement}
This study performs secondary analysis of the DAIC-WOZ corpus, which
was collected by the University of Southern California Institute for
Creative Technologies (USC ICT) under its own institutional human-
subjects protocol, with documented participant informed consent
\cite{gratch2014}. The authors did not collect any new data from human
participants and accessed only de-identified interview recordings and
questionnaire scores under the standard USC ICT data-use agreement;
no additional institutional review was required for this secondary
analysis at the authors' institution.

% ============================================================
\section*{Data Availability Statement}
The DAIC-WOZ corpus is not publicly redistributable; it is available
to qualified researchers directly from USC ICT
(\url{https://dcapswoz.ict.usc.edu/}) upon signing a data-use
agreement. FER2013, used only for encoder pre-training, is a publicly
available Kaggle dataset. All source code, trained model
configurations, and result artifacts (excluding the restricted raw
dataset) are available at
\url{https://github.com/DevMursLab/Thesis}.

\subsection*{Reproducibility Details}
All experiments use Python~3.10, PyTorch~$\geq$2.0, scikit-learn~$\geq$1.3,
and librosa~$\geq$0.10 (exact pinned versions in the repository's
\texttt{requirements.txt}). The 5-seed ablation and multi-task
experiments use seeds $\{42, 1, 7, 123, 2024\}$; the 10-seed fairness
and attention-vs-concatenation experiments use seeds
$\{42,1,7,13,21,100,2024,7777,555,88\}$. \texttt{torch.manual\_seed}
and \texttt{numpy.random.seed} are set identically per run; exact
bitwise reproducibility across GPU/CPU hardware is not guaranteed
due to non-deterministic CUDA kernels, though result variation across
reruns at fixed seed was observed to be negligible relative to the
inter-seed variance reported throughout this paper.

% ============================================================
\section*{Author Contributions}
Md.~Mursalin: conceptualization, methodology, software, formal
analysis, investigation, data curation, writing---original draft,
visualization. Md.~Shovon (Supervisor): validation, writing---review
and editing, supervision. Both authors read and approved the final
manuscript.

% ============================================================
\section*{Conflict of Interest Statement}
The authors declare that they have no known competing financial
interests or personal relationships that could have appeared to
influence the work reported in this paper.

% ============================================================
\section*{Funding Statement}
This research received no external grant funding from any
funding agency in the public, commercial, or not-for-profit
sectors.

% ============================================================
\section*{Acknowledgment}
The DAIC-WOZ corpus was provided by the University of Southern
California Institute for Creative Technologies under a restricted-data
agreement. FER2013 is a publicly available Kaggle dataset.
The authors thank the reviewers for constructive feedback.

% ============================================================
\IEEEtriggeratref{14}
\bibliographystyle{IEEEtran}

\begin{thebibliography}{19}

\bibitem{who2023}
World Health Organization, ``Depressive disorder (depression),''
\textit{WHO Fact Sheets}, Geneva, Switzerland, 2023.
[Online]. Available: \url{https://www.who.int/news-room/fact-sheets/detail/depression}

\bibitem{gratch2014}
J.~Gratch, R.~Artstein, G.~Lucas, G.~Stratou, S.~Scherer,
A.~Nazarian, R.~Wood, J.~Boberg, D.~DeVault, S.~Marsella,
D.~R.~Traum, S.~Darby, and L.-P.~Morency,
``The Distress Analysis Interview Corpus of human and computer
interviews,'' in \textit{Proc. Int. Conf. Lang. Resources Eval.
(LREC)}, Reykjavik, Iceland, 2014, pp.~3123--3128.

\bibitem{ringeval2017}
F.~Ringeval, B.~Schuller, M.~Valstar, J.~Gratch, R.~Cowie,
S.~Scherer, S.~Mozgai, N.~Cummins, M.~Schmitt, and M.~Pantic,
``AVEC 2017: Real-life depression, and affect recognition workshop
and challenge,'' in \textit{Proc. 7th ACM Int. Workshop Audio/
Visual Emotion Challenge (AVEC)}, Mountain View, CA, USA, 2017,
pp.~3--9.

\bibitem{gong2017}
Y.~Gong and C.~Poellabauer,
``Topic modeling based multi-modal depression severity assessment,''
in \textit{Proc. 7th ACM Int. Workshop Audio/Visual Emotion
Challenge (AVEC)}, 2017, pp.~53--60.

\bibitem{williamson2016}
J.~R.~Williamson, E.~Godoy, M.~Cha, A.~Schwarzenthal,
D.~Quatieri, and J.~M.~Liss,
``Detecting depression using vocal, facial and semantic
communication cues,'' in \textit{Proc. 6th Int. Workshop
Audio/Visual Emotion Challenge (AVEC)}, 2016, pp.~11--18.

\bibitem{ray2019}
A.~Ray, S.~Kumar, R.~Reddy, P.~Desai, and R.~Garg,
``Multi-level attention network using text, audio and video
for depression prediction,'' in \textit{Proc. 9th Int. Workshop
Audio/Visual Emotion Challenge (AVEC)}, 2019, pp.~81--88.

\bibitem{shen2022}
Y.~Shen, C.~Yang, and C.-Y.~Lin,
``A transformer-based multimodal depression detection system,''
\textit{IEEE Trans. Affective Comput.}, vol.~13, no.~4,
pp.~1--12, 2022.

\bibitem{covarep2014}
G.~Degottex, J.~Kane, T.~Drugman, T.~Raitio, and S.~Scherer,
``COVAREP---A collaborative voice analysis repository for speech
technologies,'' in \textit{Proc. IEEE Int. Conf. Acoust. Speech
Signal Process. (ICASSP)}, Florence, Italy, 2014, pp.~960--964.

\bibitem{scherer2015}
S.~Scherer, G.~Stratou, and L.-P.~Morency,
``Investigating voice quality as a speaker-independent indicator
of depression and PTSD,'' in \textit{Proc. Interspeech},
Dresden, Germany, 2015, pp.~847--851.

\bibitem{hardt2016}
M.~Hardt, E.~Price, and N.~Srebro,
``Equality of opportunity in supervised learning,'' in
\textit{Proc. Adv. Neural Inf. Process. Syst. (NeurIPS)},
Barcelona, Spain, 2016, pp.~3315--3323.

\bibitem{barocas2019}
S.~Barocas, M.~Hardt, and A.~Narayanan,
\textit{Fairness and Machine Learning: Limitations and Opportunities}.
2019. [Online]. Available: \url{https://fairmlbook.org}

\bibitem{ekman1978}
P.~Ekman and W.~V.~Friesen,
\textit{Facial Action Coding System (FACS): A Technique for the
Measurement of Facial Action}.
Palo Alto, CA, USA: Consulting Psychologists Press, 1978.

\bibitem{abnar2020}
S.~Abnar and W.~Zuidema,
``Quantifying attention flow in transformers,'' in
\textit{Proc. 58th Annu. Meeting Assoc. Comput. Linguistics
(ACL)}, Online, 2020, pp.~4190--4197.

\bibitem{shrikumar2017}
A.~Shrikumar, P.~Greenside, and A.~Kundaje,
``Learning important features through propagating activation
differences,'' in \textit{Proc. 34th Int. Conf. Mach. Learn.
(ICML)}, Sydney, Australia, 2017, pp.~3145--3153.

\bibitem{goodfellow2015}
I.~J.~Goodfellow, D.~Erhan, P.~L.~Carrier, A.~Courville,
M.~Mirza, B.~Hamner, W.~Cukierski, Y.~Tang, D.~Thaler,
D.-H.~Lee, Y.~Zhou, C.~Ramaiah, F.~Feng, R.~Li, X.~Wang,
D.~Athanasakis, J.~Shawe-Taylor, M.~Milakov, J.~Park,
R.~Ionescu, M.~Popescu, C.~Grozea, J.~Bergstra,
J.~Xie, L.~Romaszko, B.~Xu, Z.~Chuang, and Y.~Bengio,
``Challenges in representation learning: A report on three
machine learning contests,''
\textit{Neural Netw.}, vol.~64, pp.~59--63, Apr.~2015.

\bibitem{bouthillier2021}
X.~Bouthillier, C.~Laurent, and P.~Vincent,
``Accounting for variance in machine learning benchmarks,''
in \textit{Proc. Mach. Learn. Syst. (MLSys)}, 2021.

\bibitem{tsai2019multimodal}
Y.-H.~H.~Tsai, S.~Bai, P.~P.~Liang, J.~Z.~Kolter,
L.-P.~Morency, and R.~Salakhutdinov,
``Multimodal transformer for unaligned multimodal language
sequences,'' in \textit{Proc. 57th Annu. Meeting Assoc.
Comput. Linguistics (ACL)}, Florence, Italy, 2019,
pp.~6558--6569.

\bibitem{baltrusaitis2016}
T.~Baltru\v{s}aitis, P.~Robinson, and L.-P.~Morency,
``OpenFace: An open source facial behavior analysis toolkit,''
in \textit{Proc. IEEE Winter Conf. Appl. Comput. Vis. (WACV)},
Lake Placid, NY, USA, 2016, pp.~1--10.

\bibitem{danylenko2026}
I.~Danylenko and O.~Unold,
``Common pitfalls and recommendations for use of machine
learning in depression severity estimation: DAIC-WOZ study,''
\textit{Appl. Sci.}, vol.~16, no.~1, p.~422, 2026.

\end{thebibliography}

\end{document}
